% Figure captions for The Resonant Spectral Stack
% Each figure is a 1×4 panel (three 2-D subplots + one 3-D subplot).
% All data are computed live from the resonant_stack Python package.

% ------------------------------------------------------------------
% Panel 1
% ------------------------------------------------------------------
\begin{figure*}[!htbp]
    \centering
    \includegraphics[width=\textwidth]{figures/panel_1_dispersion.png}
    \caption{\textbf{Cauchy Dispersion and Discrete Snell Transitions across a Five-Fluid Stack.}
    (\textbf{A}) Refractive index $n(\lambda)$ for all seven modelled fluids (water, glycerol, ethanol, benzene, CS$_2$, toluene, olive oil) over the visible range 380--750\,nm, computed from the two-term Cauchy formula $n(\lambda)=A+B/\lambda^2$.  All curves decrease monotonically with wavelength (normal dispersion), spanning from $n\approx 1.33$ (water, red end) to $n\approx 1.67$ (CS$_2$, blue end).  The spread of Cauchy coefficients $(A_k,B_k)$ between layers is the sole source of the stack's spectral resolving power.
    (\textbf{B}) Ray angle inside the last fluid layer (degrees) versus wavelength for three representative incidence angles $\theta_0\in\{5^\circ,15^\circ,25^\circ\}$, after propagating through the five-layer water/glycerol/ethanol/benzene/CS$_2$ stack via the discrete Snell transition (partition depth $n=12$, angular step $\delta\theta=\pi/576\approx0.31^\circ$).  Each curve is piecewise-constant in wavelength, reflecting the quantised nature of the discrete transition; the step structure is finer at larger incidence where the dispersion-induced angular shift between wavelengths exceeds $\delta\theta$.
    (\textbf{C}) Spectral transfer matrix $A_{kj}$ (output angle in degrees) for $K=5$ wavelengths (400, 480, 550, 620, 700\,nm) and $J=5$ input angles ($5^\circ$--$25^\circ$ in $5^\circ$ steps).  No two rows are identical, confirming that the five wavelengths produce distinguishable output patterns; the matrix has numerical rank 5 (see Panel~2), satisfying the full-rank condition of the Spectral Separability Theorem.
    (\textbf{D}) Three-dimensional surface of $n(\lambda,\text{fluid})$ over wavelength (400--700\,nm) and fluid index (0--6).  The surface is a ruled manifold whose cross-sections in the wavelength direction are the individual Cauchy curves of panel~(A), and whose cross-sections in the fluid direction are the dispersion-contrast profiles $\Delta n_k(\lambda)$ that drive the rank of the transfer matrix.  The smooth, non-planar geometry of this surface is the geometric statement of Corollary~\ref{cor:spectrometer}.\label{fig:panel1_dispersion}}
\end{figure*}

% ------------------------------------------------------------------
% Panel 2
% ------------------------------------------------------------------
\begin{figure*}[!htbp]
    \centering
    \includegraphics[width=\textwidth]{figures/panel_2_separability.png}
    \caption{\textbf{Spectral Separability Theorem: Transfer Matrix Rank and Conditioning.}
    (\textbf{A}) Singular values $\sigma_i$ (degrees, log scale) of the $5\times5$ spectral transfer matrix $A_\lambda$ for the five-fluid stack.  All five singular values are strictly positive ($\sigma_1\approx90^\circ$ down to $\sigma_5\approx0.08^\circ$), spanning three decades, confirming numerical rank 5 and hence full row rank.  By Theorem~\ref{thm:separability}, this guarantees unique wavelength recovery from the observed output-angle vector.
    (\textbf{B}) Numerical rank $\mathrm{rank}(A_N)$ versus number of distinct fluid layers $N\in\{1,\ldots,5\}$, where each additional layer is chosen from a fluid not previously present in the stack.  The rank increases monotonically with $N$, saturating at the full-rank value 5 (red dashed line) when $N=5$.  Adding a fluid with a new Cauchy-coefficient pair $(A_k,B_k)$ contributes an independent row to $A_\lambda$, in agreement with the proof of Corollary~\ref{cor:spectrometer}.
    (\textbf{C}) Condition number $\kappa(A_\lambda)$ of the $5\times5$ transfer matrix versus partition depth $n$ (angular step $\delta\theta=\pi/(4n^2)$).  At shallow depths ($n\leq 8$) the condition number is astronomically large ($\kappa\sim10^{29}$) because the coarse quantisation collapses neighbouring angles onto the same discrete value, reducing the effective rank.  Above $n=10$ the condition number stabilises near $\kappa\sim10^5$, indicating that a partition depth of at least $n=10$ ($\delta\theta\approx0.14^\circ$) is required for the transfer matrix to be well-conditioned and the separability theorem to be numerically exploitable.
    (\textbf{D}) Three-dimensional surface of $A_{kj}$ (output angle in degrees) over wavelength index $k\in\{0,\ldots,4\}$ and input-angle index $j\in\{0,\ldots,4\}$.  The surface consists of five distinct shelves at approximately $5^\circ,10^\circ,15^\circ,20^\circ,25^\circ$ corresponding to the five input angles, with weak wavelength-dependent curvature within each shelf arising from the Cauchy dispersion.  The staircase geometry is the three-dimensional representation of why the rank equals the number of distinct input angles rather than the number of layers.\label{fig:panel2_separability}}
\end{figure*}

% ------------------------------------------------------------------
% Panel 3
% ------------------------------------------------------------------
\begin{figure*}[!htbp]
    \centering
    \includegraphics[width=\textwidth]{figures/panel_3_temporal.png}
    \caption{\textbf{Temporal Echo Delays: Linear Scaling with Loop Count and Shapiro Equivalence.}
    (\textbf{A}) Temporal echo delay $\Delta\tau_{n_\ell}(\lambda)$ (picoseconds) versus loop count $n_\ell\in\{1,\ldots,20\}$ for all five wavelength channels (400\,nm to 700\,nm, rainbow colour scale).  All five curves are exactly linear in $n_\ell$, confirming $\Delta\tau_{n_\ell}(\lambda)=n_\ell\,\tau_0(\lambda)$ to within floating-point precision.  The absolute delay ranges from $\approx370$\,ps (700\,nm, lower) to $\approx400$\,ps (400\,nm, upper) at $n_\ell=20$, reflecting the larger optical path length accumulated by shorter wavelengths in normally dispersive fluids.
    (\textbf{B}) Differential temporal delay $\delta\tau_{n_\ell}(400\,\text{nm},700\,\text{nm})$ versus $n_\ell$, isolating the chromatic dispersion component.  The curve is exactly linear with slope $1.617$\,ps loop$^{-1}$ (linear fit, red dashed), confirming that the spectral resolution in the time domain scales as $n_\ell\,\delta\tau_0$---each additional cavity round trip amplifies the wavelength separation by a fixed increment, analogous to increasing the number of grating rulings in a diffraction spectrometer.
    (\textbf{C}) Dispersive delay (the wavelength-dependent excess over vacuum transit) at $n_\ell=5$, comparing the optical-path-length computation (OPL, blue circles) against the Shapiro-potential formula $\Delta t_\text{Shapiro}=-2c^{-3}\int\Phi_\text{eff}\,dl$ (red dashed squares), with the relative error shown on the secondary axis (grey triangles, log scale).  The two independent calculations agree to within $10^{-8}$ across all wavelengths, validating Theorem~\ref{thm:shapiro-equiv}: the fluid-layer refractive-index profile acts as an effective gravitational potential with $-2\Phi_\text{eff}/c^2 = \mathrm{Re}[n(\lambda,z)]-1$.
    (\textbf{D}) Three-dimensional surface of $\Delta\tau_{n_\ell}(\lambda)$ over wavelength (400--700\,nm) and loop count (1--20), constituting the $(\lambda,n_\ell)$ slice of the four-dimensional output tensor $\mathcal{O}[\lambda,n_\ell,z,\text{obs}]$ at fixed depth $z$ and Beer-Lambert observation type.  The surface is a ruled bilinear manifold---planar in $n_\ell$ (linearity proven analytically) and slightly curved in $\lambda$ (Cauchy dispersion)---whose gradient vector field encodes the full chromatic-dispersion profile of the stack.\label{fig:panel3_temporal}}
\end{figure*}

% ------------------------------------------------------------------
% Panel 4
% ------------------------------------------------------------------
\begin{figure*}[!htbp]
    \centering
    \includegraphics[width=\textwidth]{figures/panel_4_triple.png}
    \caption{\textbf{Triple Observation Identity: Beer-Lambert, Chromatographic, and Kirchhoff Laws as One.}
    (\textbf{A}) Partition fraction versus optical depth $\alpha L$ (log scale, range $10^{-3}$--$10^{1}$) for the three arms of the Triple Observation Identity: transmittance $T=e^{-\alpha L}$ (Beer-Lambert, blue solid), the mobile-phase fraction $1/(1+k')$ with $k'=e^{\alpha L}-1$ (chromatographic, red dashed), and the normalised transmitted voltage $V_\text{frac}=T$ (Kirchhoff, green dotted).  The three curves are numerically indistinguishable, forming a single sigmoidal transition from 1 (optically thin, $\alpha L\ll1$) to 0 (optically thick, $\alpha L\gg1$).  This is the visual statement of Theorem~\ref{thm:triple-identity}: the three laws are related by bijective changes of variable.
    (\textbf{B}) Recovered absorption coefficient $\alpha_\text{rec}$ (log scale) versus true value $\alpha_\text{true}$ for each of the three inversion arms: Beer-Lambert (blue), chromatographic (red dashed), and Kirchhoff (green dotted).  All three lie exactly on the identity line (black, slope 1) across four decades from $10^{-4}$ to $10^{0}\,\mu\text{m}^{-1}$, confirming that each arm inverts to the same $\alpha$ without systematic bias.
    (\textbf{C}) Relative recovery error $|\alpha_\text{rec}-\alpha_\text{true}|/\alpha_\text{true}$ versus $\alpha$ for all three arms.  The Beer-Lambert and Kirchhoff arms achieve errors at the level of machine precision ($\sim10^{-15}$, well below the $10^{-12}$ validation threshold shown by the black dashed line).  The chromatographic arm errors are slightly larger due to the $\log(1+k')$ inversion but remain below $10^{-12}$ throughout the tested range.  The cross-validation spread between arms (not plotted separately) is identically zero to machine precision for the BL-Kirchhoff pair, confirming the algebraic identity $\alpha_\text{BL}\equiv\alpha_\text{KV}$.
    (\textbf{D}) Three-dimensional surface of $\log_{10}(\alpha_\text{rec})$ over $\log_{10}(\alpha L)$ (horizontal axis, range $-2$ to $1$) and arm index (Beer-Lambert, Chromatographic, Kirchhoff).  The surface is a perfectly flat plane with gradient exactly $+1$ in the $\log_{10}(\alpha L)$ direction and gradient 0 in the arm-index direction, the geometric representation that all three arms recover identical $\alpha$ at all optical depths.  The colour gradient (yellow at high $\alpha$, blue at low $\alpha$) encodes the magnitude.\label{fig:panel4_triple}}
\end{figure*}

% ------------------------------------------------------------------
% Panel 5
% ------------------------------------------------------------------
\begin{figure*}[!htbp]
    \centering
    \includegraphics[width=\textwidth]{figures/panel_5_composition.png}
    \caption{\textbf{Composition Inflation, Residue Decomposition, and the Four-Dimensional Output Tensor.}
    (\textbf{A}) Composition inflation $T(n,d)$ (semilog scale) versus partition depth / loop count $n\in\{1,\ldots,15\}$ for $d=3$ (black circles, growth rate $4^{n-1}$) and $d=2$ (grey squares, growth rate $3^{n-1}$).  Both sequences grow exponentially, confirming $T(n,d)=d(d+1)^{n-1}$ (Theorem~\ref{thm:composition-inflation}).  The dotted reference line marks $b^d=27$ (one complete base-3 partition in three-dimensional space), the denominator of the residue fraction $f_R=26/27$.  At $n=10$ the stack already supports $T(10,3)=786\,432$ distinguishable trajectory histories for $K=5$ spectral channels, providing $\log_2(786\,432/5)\approx17$ bits of redundancy per channel (Proposition~\ref{prop:redundancy}).
    (\textbf{B}) $A+R+P$ decomposition of the unit partition measure across $k=1,\ldots,20$ cycles for base $b=3$, dimension $d=3$.  The stacked area chart shows the actualised fraction $A=1/27\approx3.7\%$ (blue strip, bottom), the residue fraction $R=26/27\approx96.3\%$ (orange band, middle), and the potential $P=0$ (grey, absent) at every cycle.  The constraint $A+R+P=1$ holds exactly at each step (Theorem~\ref{thm:residue-fraction} and Definition~\ref{def:arp}).  The panel title records the derived cavity finesse $\mathcal{F}=\pi\sqrt{26/27}/(1/27)\approx83.2$ (Definition~\ref{def:finesse}), which agrees with the Fabry-P\'{e}rot finesse for a mirror reflectance equal to $f_R$.
    (\textbf{C}) Echo intensity $I_{n_\ell}(\lambda)$ (semilog scale) versus loop count $n_\ell\in\{1,\ldots,15\}$ for each of the five wavelength channels (rainbow, 400--700\,nm) and for the theoretical decay envelope $f_R^{n_\ell}$ (black dashed).  All five intensity curves decay exponentially at the rate set by the combined Beer-Lambert absorption and round-trip reflectance $f_R=26/27$; the wavelength dependence is weak because the dominant factor is $f_R^{n_\ell}$, not the per-wavelength absorption coefficient.  The close tracking of all curves to the $f_R^{n_\ell}$ envelope confirms that the cavity finesse is set by the residue fraction alone, as predicted by the partition-geometric derivation.
    (\textbf{D}) Three-dimensional surface of $\log_{10}I_{n_\ell}(\lambda)$ over wavelength (400--700\,nm) and loop count (1--15), constituting the Beer-Lambert slice $\mathcal{O}[\lambda,n_\ell,\text{BL}]$ of the full four-dimensional output tensor (Definition~\ref{def:output-tensor}).  The surface is a smoothly inclined bilinear manifold: linear in $n_\ell$ (exponential decay on the log scale) with a weak positive curvature in $\lambda$ (shorter wavelengths decay faster due to larger $\alpha(\lambda)$).  The two chromatographic and Kirchhoff slices $\mathcal{O}[\lambda,n_\ell,\text{CR}]$ and $\mathcal{O}[\lambda,n_\ell,\text{KL}]$ are related to this surface by the bijections of Theorem~\ref{thm:triple-identity} and are geometrically identical up to a change of the vertical-axis variable.\label{fig:panel5_composition}}
\end{figure*}
